\documentclass[11pt]{article}
\usepackage[margin=1in]{geometry}
\usepackage[T1]{fontenc}
\usepackage{lmodern,microtype,amsmath,amssymb,amsthm,booktabs,tabularx,enumitem,needspace}
\usepackage[colorlinks=true,linkcolor=blue,urlcolor=blue]{hyperref}
\newtheorem{proposition}{Proposition}
\newtheorem{definition}{Definition}
\newcommand{\dd}{\mathrm d}
\begin{document}
\begin{center}
{\Large Housing over the Life Cycle and Fertility}\par
\smallskip
{\large An illustrative model}\par
\smallskip
{\small Discussion draft for Tommaso De Santo \quad September 8, 2026}
\end{center}
\medskip

Young households may have the lifetime resources to buy larger homes but lack
the funds needed when they are young. We study when borrowing limits then keep
housing with old owners. The welfare comparison allows households to adjust
consumption and housing over their lives. Fertility is held fixed for this
comparison; its response to housing finance is studied separately.

\section{Environment}

\textbf{Households.} At date $t$ there are $Y_t$ young and $O_t$ old households.
A young household has liquid wealth $b_i>0$, income $y_i^y>0$ when young, and
income $y_i^o\ge0$ when old. The triple $(y_i^y,b_i,y_i^o)$ has distribution $F$
in each entering cohort. Future income is known. Households choose completed
fertility once and retain their tenure when old.

\textbf{Preferences.} Each child uses $\chi>0$ units of goods and $\kappa>0$
units of space. Total nondurable expenditure is $c$ and housing is $h$; the
bundles left for adults are:
\begin{equation}
x=c-\chi n>0,\qquad s=h-\kappa n>0.
\end{equation}
Young and old utility are:
\begin{align}
u_t^y(c,h,n)&=\log(c-\chi n)+\alpha\log(h-\kappa n)+\vartheta_t\log n,
\label{eq:new_uy}\\
u^o(c^2,h^2,e)&=\log c^2+\gamma\log h^2+\omega_B\log e.
\label{eq:new_uo}
\end{align}
All preference weights are positive. Households discount future utility by
$\beta>0$. Fertility $n>0$ is continuous; the superscript 2 labels old-age
quantities. The estate $e$ consists of financial assets and the proceeds from
selling housing at death. It enters the parent's utility but does not determine
an entering household's wealth $b_i$.

A household draws an ownership preference $\xi_i$ before making its choices.
The draw is independent of its endowments and logistic with location $\bar\xi$
and scale $\sigma_\xi>0$.

\textbf{Housing and financial markets.} The housing stock is $\bar H>0$.
Rental units satisfy $h\le h_R^{\max}$ and owner units satisfy
$h\le h_O^{\max}$, where $0<h_R^{\max}<h_O^{\max}$. Goods and bonds trade with
the rest of the world. The bond price is $q\in(0,1)$ and its gross return is
$R_f=1/q$. House prices are $P_t>0$. Owners pay property tax at rate $\tau^p$.
Competition among rental intermediaries gives:
\begin{equation}
u_t\equiv qr_t=(1+q\tau^p)P_t-qP_{t+1}>0.
\label{eq:new_usercost}
\end{equation}
Here $r_t$ is rent paid at the end of the period, and $u_t$ is its value at
the beginning. Property-tax revenue is rebated equally to young and old
households; $T_t$ denotes the rebate at the beginning of the period.

\textbf{Home finance.} Current income and liquid wealth are available at
purchase. A mortgage finances at most the share $\phi_t\in(0,1)$ of the house
price. Principal and accumulated interest are repaid on entering old age.
Unsecured borrowing is unavailable. Old owners can buy or sell housing within
the owner size limit, using their income and wealth without new borrowing.
Thus an old owner's choice is not bounded by the size of its previous home.

\textbf{Demography.} A child produces $\nu>0$ young households next period,
after survival and household formation. Cohort masses satisfy:
\begin{equation}
Y_{t+1}=\nu\bar n_tY_t,\qquad O_{t+1}=Y_t,
\label{eq:new_demography}
\end{equation}
where $\bar n_t$ is average fertility among the young. Population here counts
adult households.

\section{Household choices}

\textbf{Young renters.} Financial wealth on entering old age is $a'$.
Conditional on renting, the young household solves:
\begin{align}
W_t^R(i)=\max_{c,h,n,a'}\;&u_t^y(c,h,n)+\beta V_{t+1}^R(a';i),\nonumber\\
&c+qa'+u_th=y_i^y+b_i+T_t,\qquad a'\ge0,\quad h\le h_R^{\max}.
\label{eq:new_young_renter}
\end{align}
The renter pays for current housing services and saves for old age.

\textbf{Young owners.} For an owner, $a'$ is financial wealth net of mortgage
repayment. Its problem is:
\begin{align}
W_t^O(i)=\max_{c,h,n,a'}\;&u_t^y(c,h,n)+\beta V_{t+1}^O(a',h;i),\nonumber\\
&c+qa'+(1+q\tau^p)P_th=y_i^y+b_i+T_t,\nonumber\\
&qa'+\phi_tP_th\ge0,\qquad h\le h_O^{\max}.
\label{eq:new_young_owner}
\end{align}
To see the mortgage directly, let $d$ be principal borrowed at purchase and
$k$ the amount invested in bonds. Then $0\le d\le\phi_tP_th$, $k\ge0$,
$a'=(k-d)/q$, and the budget is
$c+k+(1+q\tau^p)P_th=y_i^y+b_i+T_t+d$.
Eliminating $k,d$ gives \eqref{eq:new_young_owner}.

\textbf{Old households.} An old owner has net financial wealth $a$ and title
to $H$ units of housing. Let $a^e\ge0$ be financial saving during old age.
The estate is:
\begin{equation}
e=\begin{cases}
q^{-1}a^e,&\text{renter},\\
q^{-1}a^e+P_{t+1}h^2,&\text{owner}.
\end{cases}
\label{eq:new_estate}
\end{equation}
The retained house is sold at death, at the end of old age. The two old-age
problems are:
\begin{align}
V_t^R(a;i)=\max_{c^2,h^2,e}\;&u^o(c^2,h^2,e),\nonumber\\
&c^2+qe+u_th^2=a+y_i^o+T_t,\quad h^2\le h_R^{\max},
\label{eq:new_old_renter}\\
V_t^O(a,H;i)=\max_{c^2,h^2,e}\;&u^o(c^2,h^2,e),\nonumber\\
&c^2+qe+u_th^2=a+P_tH+y_i^o+T_t,\nonumber\\
&h^2\le h_O^{\max},\qquad e\ge P_{t+1}h^2.
\label{eq:new_old_owner}
\end{align}
The owner's estate restriction is equivalent to $a^e\ge0$. Its cash budget
before substituting for the estate is
$c^2+a^e+(1+q\tau^p)P_th^2=a+P_tH+y_i^o+T_t$.
Income, liquid wealth, and sale proceeds therefore enter the same budget.

\textbf{Tenure.} The household owns if $W_t^O(i)+\xi_i\ge W_t^R(i)$.
Its ownership probability is:
\begin{equation}
\pi_t^O(i)=
\frac{\exp\{[W_t^O(i)+\bar\xi]/\sigma_\xi\}}
{\exp\{W_t^R(i)/\sigma_\xi\}+\exp\{[W_t^O(i)+\bar\xi]/\sigma_\xi\}}.
\label{eq:new_tenure}
\end{equation}
The taste draw affects tenure but not choices conditional on tenure.

\section{Equilibrium and financing}

An equilibrium determines house prices, rebates, household choices, and the
distribution of assets and housing across cohorts. Let $G_t$ denote the
distribution of old households over financial wealth, housing titles, old
income, and tenure. Housing and the government budget satisfy:
\begin{equation}
Y_t\bar h_t^y+O_t\bar h_t^o=\bar H,\qquad
(Y_t+O_t)T_t=q\tau^pP_t\bar H.
\label{eq:new_markets}
\end{equation}
The bars average occupied housing over types and tenure choices.

\begin{definition}
Given the initial cohorts and old distribution, an equilibrium is a sequence
of prices, rebates, household choices, tenure probabilities, and cohort
distributions such that households solve
\eqref{eq:new_young_renter}--\eqref{eq:new_old_owner}, tenure follows
\eqref{eq:new_tenure}, rental entry satisfies \eqref{eq:new_usercost}, markets
satisfy \eqref{eq:new_markets}, and cohorts evolve according to
\eqref{eq:new_demography} and the young households' choices.
\end{definition}

At a positive stationary equilibrium, prices, distributions, and cohort
masses are constant. In particular:
\begin{equation}
Y=O=N,\qquad \bar n=1/\nu,\qquad
N=\frac{\bar H}{\mathbb E[h^y+h^o]}.
\label{eq:new_stationary}
\end{equation}
The housing stock determines cohort mass once stationary household choices
are known. Appendix~\ref{app:new_existence} gives explicit conditions for
existence with heterogeneous endowments and positive property taxes.

\textbf{Current and lifetime resources.} At stationary prices write
$w=y_i^y+b_i+T$, $v=y_i^o+T$, and:
\begin{equation}
p=(1-q+q\tau^p)P,\qquad L=(1-\phi+q\tau^p)P.
\end{equation}
Here $p$ is the service cost of housing and $L$ is its cash requirement.
For an owner with old-age resources $z=a'+Ph+v$, the two budgets become:
\begin{equation}
c+ph+qz=w+qv,\qquad c+Lh\le w.
\label{eq:new_two_budgets}
\end{equation}
The first values the household's lifetime resources; the second requires
current consumption and the cash portion of housing to be paid when young.

When old housing is below its cap, the old value has the form
$(1+\gamma+\omega_B)\log z+$constant, even if its estate floor binds.
Write $K=1+\gamma+\omega_B$, $B=\beta K$, and
$D=1+\alpha+\vartheta+B$. If the young and old caps are slack at the
allocation without the mortgage restriction, that allocation has:
\begin{equation}
\begin{aligned}
c^*&=\frac{w+qv}{D}
       \left(1+\frac{\vartheta\chi}{\chi+p\kappa}\right),\\
h^*&=\frac{w+qv}{D}
       \left(\frac\alpha p+\frac{\vartheta\kappa}{\chi+p\kappa}\right).
\end{aligned}
\label{eq:new_free_choices}
\end{equation}
The mortgage limit strictly restricts the household precisely when
$c^*+Lh^*>w$. Equivalently:
\begin{equation}
\boxed{\quad
\frac{w+qv}{w}>
\frac D{1+\alpha L/p+
\vartheta(\chi+L\kappa)/(\chi+p\kappa)}.
\quad}
\label{eq:new_binding}
\end{equation}
The household wants an allocation whose current consumption and down payment
exceed current resources. If a physical cap binds, the same cash test applies
to the allocation computed with that cap retained.

\section{Housing allocation and welfare}

We ask whether financing limits prevent a mutually beneficial reallocation
of housing over the life cycle. The planner may relax these financing limits
and make household-specific transfers and loans with future repayment.
It respects the housing stock, tenure size limits,
existing creditors' claims, and the external bond return. Fertility and tenure
are held fixed. Selected households may change their own old-age consumption,
housing, and estates; other households must be compensated.

\begin{proposition}[Housing over the life cycle]
Consider a strictly mortgage-constrained young owner whose young and old
housing caps and old estate floor are slack. If
\begin{equation}
\boxed{\quad
\beta(1+\gamma+\omega_B)(1-\phi+q\tau^p)\ge\phi-q,
\quad}
\label{eq:new_housing_condition}
\end{equation}
relaxing its financing restriction, at fixed prices and fixed fertility,
increases its desired young housing and reduces its old-age housing. A finite
step in this direction is feasible below the physical owner cap and strictly
improves lifetime utility.
\label{prop:new_lifetime}
\end{proposition}

The condition permits $q<\phi$. Greater weight on old-age utility makes it
easier to satisfy. This result concerns the household's allocation over its
life: young housing need not initially have a higher marginal value in units
of current consumption than old housing.

\begin{proof}
Hold equilibrium fertility $n$ fixed. Set
$w_n=w-(\chi+L\kappa)n$ and $M_n=w+qv-(\chi+p\kappa)n$.
The adult allocation solves:
\begin{equation}
\max\ \log x+\alpha\log s+B\log S,\qquad
x+ps+S=M_n,\quad x+Ls\le w_n,
\label{eq:new_fixed_n}
\end{equation}
where $S=qz$. Without the financial restriction its choice is
$(x_F,s_F,S_F)=M_n(1,\alpha/p,B)/(1+\alpha+B)$.

Let $\lambda,\mu>0$ be the lifetime-budget and cash multipliers at the
constrained allocation. Define $r=\mu/(\lambda+\mu)$,
$\rho=(1-r)p+rL$, and $j=1/(\lambda+\mu)$. The first-order conditions give
$(x_C,s_C,S_C)=(j,\alpha j/\rho,Bj/(1-r))$. Subtraction yields:
\begin{equation}
s_C-s_F=
\frac{\alpha jr\{(p-L)(1-r)-BL\}}
{p(1+\alpha+B)\rho(1-r)}.
\label{eq:new_housing_identity}
\end{equation}
Condition \eqref{eq:new_housing_condition} is $BL\ge p-L$ and makes this
expression strictly negative. Also:
\begin{equation}
\frac{S_C}{M_n}
=\frac B{B+(1-r)(1+\alpha p/\rho)}
>\frac B{1+\alpha+B}=\frac{S_F}{M_n}.
\end{equation}
Old resources, and therefore old housing, fall. A sufficiently small convex
step toward the unrestricted allocation preserves the young cap, the old
cap, and the estate floor while strictly increasing utility.
\end{proof}

The inequality is sufficient rather than necessary. When $p>L$, an exact
condition for a strictly constrained household to prefer larger young housing
in the fixed-fertility comparison is:
\begin{equation}
Lqv>(p-L)(w-\chi n).
\label{eq:new_exact_income}
\end{equation}
Enough future income relative to current resources also delivers the direction
when \eqref{eq:new_housing_condition} fails.

\begin{proposition}[Inefficiency]
Suppose the stationary equilibrium has a positive share of young owners
satisfying Proposition~\ref{prop:new_lifetime}. It admits a Pareto improvement
that transfers housing from old owners to young owners at the intervention
date. Fertility, tenure, and every future cohort size are unchanged. New
external borrowing is repaid by the selected households' estate date.
\label{prop:new_pareto}
\end{proposition}

The required owner share follows from explicit income, preference, and
capacity inequalities in Appendix~\ref{app:new_existence}; it is not an
assumed difference in marginal housing values. Under those inequalities a
stationary equilibrium exists, and every stationary equilibrium satisfies the
proposition's conditions.

The argument uses two housing reallocations. Current old owners supply the
additional homes needed by the selected young. When those young households
become old, their smaller homes release space to other old owners. The latter
groups are compensated exactly. A small portion of the selected households'
lifetime gain covers the compensation costs. Appendix~\ref{app:new_pareto}
gives the transfers and loan repayment.

This establishes one Pareto improvement with the stated housing direction.
It does not characterize the entire first-best allocation or establish that
a uniform mortgage reform is Pareto improving.

\section{Fertility and ownership}

For comparisons at a given transition date, the formulas in this section use
$p=u_t$ and $L=(1-\phi_t+q\tau^p)P_t$, with $v=y_i^o+T_{t+1}$.
The household chooses children by comparing their benefit with their goods
and space costs:
\begin{equation}
\frac{\vartheta}{n}
=\frac\chi{c-\chi n}+\frac{\alpha\kappa}{h-\kappa n}.
\label{eq:new_fertility}
\end{equation}
More housing lowers the space cost of another child at fixed consumption.
Along a change in both $c$ and $h$, differentiation gives:
\begin{equation}
\dd n=
\frac{(\chi/x^2)\dd c+(\alpha\kappa/s^2)\dd h}
{\vartheta/n^2+\chi^2/x^2+\alpha\kappa^2/s^2}.
\label{eq:new_fertility_differential}
\end{equation}
The financing of additional housing therefore matters for fertility.

\begin{proposition}[Fertility under easier mortgage finance]
Hold the price and rebate path fixed, and suppose old housing caps are slack
at the relevant conditional choices. A strictly constrained owner increases
fertility when $\phi$ rises if $L\ge p$. When $L<p$, a sufficient condition is:
\begin{equation}
\boxed{\quad
Bp\left(L+\frac\chi\kappa\right)
\ge(p-L)\left(p-\alpha\frac\chi\kappa\right).
\quad}
\label{eq:new_fertility_condition}
\end{equation}
With slack young housing, this condition is sharp for a fertility increase
uniform over current and future incomes at given $L,p$.
\label{prop:new_fertility}
\end{proposition}

This restriction compares the goods and space costs of children and the
weight on future utility. It contains no borrowing multiplier. The simpler
condition $(\alpha+B)\chi\ge p\kappa$ guarantees the fertility direction
throughout any range of nonnegative $L$. A binding young housing cap fixes
housing; easier finance then raises fertility by releasing current resources.

Tenure changes require one additional comparison. A higher $\phi$ expands
the owner's feasible set, so ownership weakly increases at fixed prices. For
each type, mean fertility changes according to:
\begin{equation}
\frac{\dd\bar n_i}{\dd\phi}
=\pi_i^O\frac{\dd n_i^O}{\dd\phi}
 +\frac{\dd\pi_i^O}{\dd\phi}(n_i^O-n_i^R).
\label{eq:new_selection}
\end{equation}
The fertility difference between owners and renters determines the contribution
of households switching tenure.

\begin{proposition}[Fertility including tenure choice]
Suppose $L\le p$, old housing caps are slack at both conditional choices, and:
\begin{equation}
\frac p{\alpha+B}\le\frac\chi\kappa\le\frac p\alpha.
\label{eq:new_fertility_interval}
\end{equation}
Then $n_i^O\ge n_i^R$, and easier mortgage finance weakly increases fertility
averaged over a fixed heterogeneous cohort, including tenure changes. The
result extends to finite reforms for which the conditions hold throughout.
\end{proposition}

Renting is the same current-resource problem with cash cost $L=p$ and the
smaller rental cap. Under the upper bound in
\eqref{eq:new_fertility_interval}, access to larger housing at $L=p$ weakly
raises fertility. Under the lower bound, lowering $L$ to the owner's cash
cost weakly raises it further. This gives the tenure ordering; both terms in
\eqref{eq:new_selection} are then nonnegative. The proof and the sharp
individual condition are given in Appendix~\ref{app:new_fertility}.

These are household and fixed-cohort results. The price $p$ is endogenous in
equilibrium. A mortgage reform can change prices and rebates as well as the
household's financing opportunities; the displayed restrictions do not by
themselves sign that general-equilibrium response.

\section{A fertility decline and a later housing policy}

The intended experiment starts from a stationary equilibrium. A persistent
decline in $\vartheta_t$ initiates demographic adjustment. At a later date a
housing policy is introduced, and its path is compared with continued
adjustment without the policy. The two paths share the inherited state and
non-policy shocks at announcement. If policy is anticipated, expectations
must enter choices from that earlier date.

At fixed prices and rebates, a decline in $\vartheta$ weakly reduces fertility,
including tenure selection. For each type, the envelope theorem gives
$\partial(W^O-W^R)/\partial\vartheta=\log(n^O/n^R)$. Consequently:
\begin{equation}
\begin{aligned}
\frac{\partial\bar n_i}{\partial\vartheta}
={}&\pi_i^O\frac{\partial n_i^O}{\partial\vartheta}
 +(1-\pi_i^O)\frac{\partial n_i^R}{\partial\vartheta}\\
&+\frac{\pi_i^O(1-\pi_i^O)}{\sigma_\xi}
 (n_i^O-n_i^R)\log(n_i^O/n_i^R)\ge0.
\end{aligned}
\label{eq:new_theta_selection}
\end{equation}
Conditional fertility is nondecreasing in its preference weight, and the
selection term is nonnegative regardless of which tenure has higher fertility.

For any two existing equilibrium paths with the same young cohort at date
$t_0$, demographic accounting implies:
\begin{equation}
\frac{Y_t^{\rm policy}}{Y_t^{\rm baseline}}
=\prod_{s=t_0}^{t-1}
 \frac{\bar n_s^{\rm policy}}{\bar n_s^{\rm baseline}}.
\label{eq:new_population_product}
\end{equation}
This relates population differences to fertility differences accumulated
along the path. At positive stationary endpoints, fertility returns to
$1/\nu$ in both economies; their population levels can differ.

At those endpoints, the population comparison is:
\begin{equation}
\frac{Y^{\rm policy}+O^{\rm policy}}
{Y^{\rm baseline}+O^{\rm baseline}}
=\frac{\mathbb E[h^y+h^o]^{\rm baseline}}
{\mathbb E[h^y+h^o]^{\rm policy}}.
\label{eq:new_population_endpoint}
\end{equation}
With fixed total housing, a larger stationary population requires less housing
per household over the life cycle, after prices and tenure have adjusted.

A transition existence or sign theorem requires more than stationary
existence. At every date the housing caps require
$Y_t+O_t\ge\bar H/h_O^{\max}$, and an inherited owner must have
$a+P_tH+y_i^o+T_t>0$ to finance positive old-age consumption and an estate.
The current results do not establish an invariant region satisfying these
requirements or convergence after a finite preference decline. Equations
\eqref{eq:new_population_product}--\eqref{eq:new_population_endpoint} are exact
conditional comparisons, not a proof that a mortgage policy raises fertility
at every transition date or the final population. The earlier model's
transition construction cannot be transferred to this specification unchanged.

\appendix
\section{Primitive conditions and stationary existence}
\label{app:new_existence}

This appendix supplies conditions on endowments and parameters that ensure
the stationary equilibrium and owner groups used in the welfare proposition.
They are sufficient and conservative. No price or borrowing multiplier is an
input to these conditions.

Assume bounded endowments, $w_{0i}=y_i^y+b_i\ge\underline w>0$ and
$v_{0i}=y_i^o\ge0$, and $0\le\tau^p<2$. Define:
\begin{equation}
\begin{gathered}
d_p=1-q+q\tau^p,\quad d_L=1-\phi+q\tau^p,\quad
d_{\max}=\max\{d_p,d_L\},\\
\bar M_0=\mathbb E[w_0+qv_0],\qquad
\bar T=\frac{\tau^p\bar M_0}{(1-q)(2-\tau^p)},\qquad
\bar M=\bar M_0+(1+q)\bar T.
\end{gathered}
\end{equation}
The scalar $\bar T$ bounds a balanced stationary rebate. Suppose:
\begin{equation}
h_R^{\max}>\frac\kappa\nu,\qquad
\vartheta\nu>\frac{\chi D}{\underline w}
 +\frac{\alpha\kappa}{h_R^{\max}-\kappa/\nu}.
\label{eq:new_existence_conditions}
\end{equation}
These conditions ensure that sufficiently inexpensive housing permits
fertility above replacement. They do not require renters to have fertility
above replacement at the equilibrium price.

At least one positive stationary equilibrium exists, and every stationary
equilibrium satisfies:
\begin{equation}
P_-<P<P_+,\qquad 0\le T\le\bar T,\qquad
P_-\equiv\frac{\alpha\underline w}{D d_{\max}h_R^{\max}},\quad
P_+\equiv\frac{\nu\bar M}{d_p\kappa}.
\label{eq:new_price_bounds}
\end{equation}
To establish this, combine the two age budgets. In either tenure:
\begin{equation}
x+ps+(\chi+p\kappa)n+qc^2+qp h^2+q^2e=w+qv.
\label{eq:new_lifetime_full}
\end{equation}
Scaling the first-order conditions, including both caps and the homogeneous
estate floor, gives $\lambda(w+qv)+\mu w\le D$, while
$1/x=\lambda+\mu$. Thus $x\ge w/D\ge\underline w/D$.
If the young housing cap is slack, $\alpha x/s\le d_{\max}P$. If it binds,
housing is at least $h_R^{\max}$. At $P\le P_-$ every conditional choice
therefore has $h\ge h_R^{\max}$. The fertility condition
\eqref{eq:new_fertility} and \eqref{eq:new_existence_conditions} then rule out
$n\le1/\nu$.

Conversely, \eqref{eq:new_lifetime_full} implies
$\bar n<\{\bar M_0+(1+q)T\}/(\chi+d_pP\kappa)$, so fertility is below
replacement for $P\ge P_+$ and $T\le\bar T$. Also, if
$\bar h=\mathbb E[h^y+h^o]$, the balanced-rebate map satisfies:
\begin{equation}
\mathcal T(P,T)=\frac{q\tau^pP\bar h(P,T)}2
\le\frac{\tau^p}{2d_p}\{\bar M_0+(1+q)T\}\le\bar T.
\end{equation}
The last inequality follows from
$2d_p-\tau^p(1+q)=(1-q)(2-\tau^p)>0$.

Conditional allocations are unique and continuous, including where caps
change status. Logistic tenure probabilities preserve aggregate continuity.
Map $T$ to $\mathcal T(P,T)$ and move $P$ in the direction of
$\nu\bar n(P,T)-1$, restricting the new price to $[P_-,P_+]$.
Brouwer's theorem gives a fixed point. The strict fertility signs exclude
the price boundaries, so it has replacement fertility. Setting
$N=\bar H/\bar h$ closes housing and the fiscal budget, and the young choices
generate the stationary old distribution. The same bounds apply to any
stationary equilibrium. Uniqueness and stability are separate questions.

\textbf{Eligible young owners.} Put $\ell=d_L/d_p$ and:
\begin{equation}
C_{\min}=1+\alpha\ell+\vartheta\min\{1,\ell\},\qquad
k=\frac D{C_{\min}}-1.
\end{equation}
Choose a positive-$F$-mass set $\mathcal S$ of endowments satisfying:
\begin{equation}
qv_{0i}>kw_{0i}+\max\{k-q,0\}\bar T\quad(i\in\mathcal S).
\label{eq:new_primitive_income}
\end{equation}
This ensures enough future income relative to current resources for finance
to bind, even after the common rebate. Let
$M_{\mathcal S}\ge\sup_{i\in\mathcal S}
\{w_{0i}+qv_{0i}+(1+q)\bar T\}$ and impose:
\begin{equation}
\omega_Bd_p>q\gamma,\qquad
h_O^{\max}>
\frac{M_{\mathcal S}}{d_pP_-}
\max\left\{1,\frac\gamma{qK}\right\}.
\label{eq:new_primitive_caps}
\end{equation}
The first condition makes the old estate floor slack. The lifetime budget
bounds young housing by $M/p$ and old resources by $M/q$, so the second makes
both caps slack in the actual and finance-relaxed choices of these types.
This cap bound uses the renter cap in $P_-$; enlarging the owner cap does not
lower the price bound.

The bracket in the cash-demand test \eqref{eq:new_binding} is at least
$C_{\min}$. Condition \eqref{eq:new_primitive_income} implies
$(w+qv)/w>D/C_{\min}$ throughout $T\in[0,\bar T]$.
Thus finance is strictly restrictive for every type in $\mathcal S$ at every
stationary equilibrium. The logistic distribution assigns each type a
positive owner probability. Their predecessors supply current old owners,
and unselected contemporaries supply next period's old recipients. There is
no separate equal-mass assumption on young and old income groups.

These conditions, together with \eqref{eq:new_housing_condition}, imply
Proposition~\ref{prop:new_pareto}. They establish an equilibrium result in
primitives. Their conservative bounds should not be read as evidence that
they fit a particular calibration.

\section{Compensation and external repayment}
\label{app:new_pareto}

This appendix gives the construction first for common conditional bundles;
the same bounds apply uniformly on positive-mass subsets of heterogeneous
types. Hold every selected household's fertility and tenure fixed. Take a
finite improving step from Proposition~\ref{prop:new_lifetime}. Write its
changes as $(\Delta c,\Delta h,\Delta c^2,\Delta h^2,\Delta e)$, with
$\Delta h>0$ and $\Delta h^2<0$. The lifetime budget gives:
\begin{equation}
\Delta c+p\Delta h+q\Delta c^2+qp\Delta h^2+q^2\Delta e=0.
\label{eq:new_delta_budget}
\end{equation}
If the step gives utility gain $g>0$ and new adult consumption $x_A>0$,
deduct $\eta=x_A(1-e^{-g/2})$. This leaves gain $g/2$ and releases $\eta$
goods per selected household.

Choose a small positive mass $\epsilon$ of such households. Current old
donors, of mass $m_D>0$, each surrender
$s_D=\epsilon\Delta h/m_D$ housing. Hold their estates fixed and give each:
\begin{equation}
D(s_D)=c_D\{(1-s_D/h_D)^{-\gamma}-1\}.
\end{equation}
This preserves old utility exactly. Next period give
$s_R=\epsilon(-\Delta h^2)/m_R$ extra housing to each member of an unselected
old owner group of mass $m_R>0$. Keeping their estates fixed, collect:
\begin{equation}
R(s_R)=c_R\{1-(1+s_R/h_R)^{-\gamma}\}.
\end{equation}
Their lifetime utilities are also unchanged. Small $\epsilon$ preserves their
old size caps and positive financial estates. Housing clears at both dates,
within the owner segment.

Both old groups have $\gamma c/h=p$ before the transfer. Their total
compensation costs therefore equal $p\epsilon\Delta h$ at the first date
and their total payments equal $p\epsilon(-\Delta h^2)$ at the next, up to
terms bounded by $C\epsilon^2$ for a finite constant $C$.
The net present goods requirement is at most
$C\epsilon^2-\epsilon\eta<0$ for sufficiently small positive $\epsilon$.
This is an analytical bound from the second derivatives of the two displayed
compensation functions.

In detail, the additional goods uses at the three affected dates are:
\begin{equation}
\begin{aligned}
G_0&=\epsilon(\Delta c-\eta)+m_DD(s_D),\\
G_1&=\epsilon\Delta c^2-m_RR(s_R),\\
G_2&=\epsilon\Delta e.
\end{aligned}
\end{equation}
The terminal balance of new external borrowing is
$G_0/q^2+G_1/q+G_2\le0$. Thus every new loan is repaid by the selected
households' estate date. Internal housing title payments cancel. Existing
mortgage claims and the other households' estates are preserved. Selected
households may leave smaller estates as part of their preferred lifetime
allocation; that utility cost is included in $g$.

For heterogeneous groups, first reserve a positive-mass group of future old
recipients, and then select reform households from the remaining eligible
owners. Restrict to positive-mass subsets with uniform strict margins and
$\eta_i\ge\underline\eta>0$, $\Delta h_i\le\overline H$,
$-\Delta h_i^2\le\overline K$. These subsets exist by taking countable level
sets of the positive gains and margins. The continuous ownership taste allows
any sufficiently small selected mass, even if the income distribution has
atoms. Define total housing changes by:
\begin{equation}
H_\epsilon=\int_{\mathcal S_\epsilon}\Delta h_i\,\dd\mu,
\qquad K_\epsilon=-\int_{\mathcal S_\epsilon}\Delta h_i^2\,\dd\mu.
\end{equation}
Take $H_\epsilon/m_D$ from each current donor and give $K_\epsilon/m_R$ to
each future recipient, applying that household's own compensation function.
If $\overline A_D,\overline A_R$ bound their second derivatives, the cost
above the linear housing value is at most:
\begin{equation}
\frac{\overline A_DH_\epsilon^2}{2m_D}
+q\frac{\overline A_RK_\epsilon^2}{2m_R}
\le C\epsilon^2.
\end{equation}
Fees yield at least $\underline\eta\epsilon$. Choose selected mass below
$\underline\eta/(2C)$ and small enough to preserve the uniform donor and
recipient margins. This establishes the same strict resource surplus with
unequal group masses and heterogeneous individual bundles.

\textbf{Individual budgets.} Aggregate resources can be implemented while
preserving each private mortgage. Write $A=(1+q\tau^p)P=p+qP$.
For selected household $i$, hold its original private net financial claim
$a'_{C,i}$ fixed and give transfers:
\begin{equation}
f_{0i}=\Delta c_i-\eta_i+A\Delta h_i,\qquad
f_{1i}=\Delta z_i-P\Delta h_i,
\qquad f_{0i}+qf_{1i}=-\eta_i.
\end{equation}
This is a loan of $q(P\Delta h_i-\Delta z_i)>0$ when young, repaid on
entering old age, together with fee $\eta_i$. After repaying, its old resources
are exactly $z_{C,i}+\Delta z_i>0$; its new old financial saving is
$q(e_{A,i}-Ph_{A,i}^2)>0$. No new borrowing in old age is needed.

A donor who gives up $s_D$ receives the additional transfer
$f_{Di}=D_i(s_D)-ps_D$ and raises old financial saving by $qPs_D$, keeping
its estate fixed. A recipient who obtains $s_R$ receives
$f_{Ri}=ps_R-R_i(s_R)$ and reduces financial saving by $qPs_R$, within its
strict margin. The program's present-value cost is:
\begin{equation}
\int_{\mathcal S_\epsilon}(f_{0i}+qf_{1i})\,\dd\mu
+\int_Df_{Di}\,\dd\mu+q\int_Rf_{Ri}\,\dd\mu
\le C\epsilon^2-\underline\eta\epsilon<0.
\end{equation}
The program can settle its account at the next date; the selected old
households' portfolios finish at their estate date. At the first date extra
young titles equal current donors' title reductions. At the next date those
inherited titles offset the donors' smaller death sales, and selected old
housing reductions equal recipients' additions. Their title changes at death
therefore cancel. Future entrants can make their original purchases, and all
originally contracted creditors receive their promised payments.

\section{Fertility signs}
\label{app:new_fertility}

Write $\eta=\chi/\kappa$, rescale children as $k=\kappa n$, and let $S=qz$.
At a strictly constrained owner optimum with slack caps, put
$t=Bx/S\in(0,1)$ and $\rho=\alpha x/s=(1-t)L+tp$.
These quantities are used in the proof, not in the statement of the
fertility condition. Define $H(\rho)=\alpha/\rho+\vartheta/(\eta+\rho)$
and $A=\rho^2/\alpha$. Differentiating the two choice equations for $h,k$
shows that the fertility derivative has the sign of:
\begin{equation}
\mathcal N=(1-t)(A-L\eta)+H(\rho)
\left[A(L+\eta)+\frac{t^2(L-p)}B(A-p\eta)\right].
\label{eq:new_N}
\end{equation}
The Hessian determinant is positive, by strict concavity.

The expression excluding its last term is bounded below by
$(1-t)A+\rho L+\eta tp>0$. For $L\ge p$, if the last term is negative,
$S=qv+(L-p)h$ and $v\ge0$ give $t(L-p)H(\rho)\le B$.
This bounds the entire expression below by $A+\rho L>0$.

For $L<p$, the last term is nonnegative when $A\le p\eta$. Otherwise the
first term is positive and
$t^2(1-p\eta/A)<1-\alpha\eta/p$.
Condition \eqref{eq:new_fertility_condition} makes the remaining bracket
positive. As $t$ approaches one from below, its limiting sign is the sign of:
\begin{equation}
\frac{p^2}{\alpha}(L+\eta)
-\frac{p-L}{B}\left(\frac{p^2}{\alpha}-p\eta\right).
\end{equation}
If the stated condition fails, this expression is negative, and nearby
strictly constrained income profiles give decreasing fertility. Such profiles
are feasible: choose any $x>0$ and set
$w=x\{1+\alpha+\vartheta+t(L-p)H(\rho)\}$ and
$qv=x\{B/t-(L-p)H(\rho)\}$.

If the young cap binds, differentiation at fixed housing gives:
\begin{equation}
\frac{\partial n}{\partial\phi}
=\frac{P\chi h/x^2}
{\vartheta/n^2+\chi^2/x^2+\alpha\kappa^2/s^2}>0.
\end{equation}
Continuity at changes of constraint status extends the weak sign to finite
reforms that satisfy the condition throughout.

For the tenure comparison, first set $L=p$ and raise the young cap from
$h_R^{\max}$ to $h_O^{\max}$. If finance binds, the sign of the fertility
response is the sign of $\alpha\kappa/s^2-\chi p/x^2$. The cap condition
$\alpha x/s\ge p$ and $\chi/\kappa\le p/\alpha$ make it nonnegative.
If finance is slack, the same argument replaces $\chi p$ by
$\chi p/(1+B)$ and is also nonnegative. Lowering $L$ to its actual owner
value then raises fertility under $(\alpha+B)\chi\ge p\kappa$.
Old conditional-value constants affect tenure probabilities but not these
conditional fertility comparisons.

\section{Repayment timing and the existing note}

The proposal changes three features of the existing note: current income is
available at purchase, households have an old-age income endowment, and old
owners can resize their homes. It keeps the original mortgage repayment
convention. The welfare comparison permits changes to selected households'
own future allocations; this is essential for
Proposition~\ref{prop:new_lifetime} when $q<\phi$.

For literature context, van Doornik, Fazio, Ramadorai, and Skrastins's
\href{https://bcb.gov.br/content/publicacoes/WorkingPaperSeries/WP612.pdf}{\emph{Housing and Fertility}}
(December 2024, Section 2) provides a simple
goods-and-space account of fertility but does not model endogenous tenure
choice or a mortgage down payment. Coven, Golder, Gupta, and Ndiaye's
\href{https://abdouecon.github.io/research/papers/Property_Tax.pdf}{\emph{Property Taxes and Housing Allocation Under Financial Constraints}}
(August 2026, Section 3.1) derives capitalization and intergenerational
redistribution in a simple collateral model. Its old households sell housing;
it does not supply the old-versus-young housing-services welfare result above.
These are related ingredients, not proofs of this model's conclusions.

\end{document}
